\documentclass[11pt,titlepage]{ltjsbook}
%二段組にするなら[, twocolumn]

% =========================
% 数学
% =========================
\usepackage{amsmath}
\usepackage{amssymb}
\usepackage{mathtools}
\usepackage{bm}
\usepackage{cases}
\usepackage{physics}
\usepackage{siunitx}

% =========================
% 図・色
% =========================
\usepackage{graphicx}
\usepackage{xcolor}
\usepackage{float}
\usepackage{subcaption}

% =========================
% TikZ
% =========================
\usepackage{tikz}
\usepackage{tikz-3dplot}
\usetikzlibrary{math}
\usetikzlibrary{arrows.meta}

% =========================
% 枠・装飾
% =========================
\usepackage[framemethod=tikz]{mdframed}
\usepackage{fancybox}
\usepackage{ascmac}

% =========================
% レイアウト
% =========================
\usepackage[top=20mm, bottom=20mm, left=20mm, right=20mm]{geometry}
\usepackage{nidanfloat}

% =========================
% 日本語
% =========================
\usepackage{luatexja}
\usepackage{luatexja-fontspec}
%\usepackage[deluxe]{otf}

% =========================
% コード表示
% =========================
\usepackage{listings}
\usepackage{jvlisting}

% =========================
% その他
% =========================
\usepackage{url}
\usepackage{hyperref}
\usepackage{appendix}


% =========================
% タイトル
% =========================
\title{\huge{同時測定}}
\author{\vspace{3cm}\\
\Large{担当教員}\\
\vspace{2cm}\Large{増渕達也先生}\\
\Large{04B24032}\\
\Large{笹伶夷}\\
\vspace{2cm}\Large{共同実験者：前中樹}\\
\Large{実験日：2026年4月22日、4月28日、5月12日}\\
\Large{5月13日、5月19日、5月20日、5月26日}
}
\date{}

\begin{document}
\maketitle
\tableofcontents
\clearpage


\chapter{はじめに}
\chapter{実験手法}

\chapter{実験1（電子・陽電子対消滅の観測）}
\section{実験の目的}
電子・陽電子対消滅を2つの検出器を用いた同時計測によって観測する。また、本実験によりエネルギー・時間スペクトル測定についての理解を深める。


\chapter{実験2}

\appendix
\chapter{$\gamma$線の同時測定の角度分布の導出}
\section{問題設定}
本章では、実験1で行った$\gamma$線の同時測定の角度分布の導出を与える。\par
\begin{tikzpicture}
    \draw[->, thick, >=Latex]  (-7.5, 0) -- (7.5, 0);
    \draw[->, thick, >=Latex]  (0, -5) -- (0, 5);
    \draw[thick] (-0.25, 1) -- (0.25, 1);
    \draw[thick] (-0.25, -1) -- (0.25, -1);
    \draw[thick] (0.25, 1) -- (0.25, -1);
    \draw[thick] (-0.25, -1) -- (-0.25, 1);
    \draw[arrows = {-Stealth[]}]  (-0.3, -1.2) -- (0.3, -1.2);
    \draw[arrows = {-Stealth[]}]  (0.3, -1.2) -- (-0.3, -1.2);
    \draw[thick] (2, 2) -- (2, -2);
    \draw[thick] (2, 2) -- (6, 2);
    \draw[thick] (2, -2) -- (6, -2);
    \draw[arrows = {-Stealth[]}]  (1.8, 0.1) -- (1.8, 2);
    \draw[arrows = {-Stealth[]}]  (1.8, 2) -- (1.8, 0.1);



    \node[above] at (7.5, 0) {$x$};
    \node[left] at (0, 5) {$y$};
    \node at (0.2,-1.5) {$d$};
    \node[below left] at (2, 0) {$c$}; 
    \node[left] at (1.8, 1) {$R_b$};


\end{tikzpicture}

\end{document}